%% ma: L11-B16-0047
%% lop: 11
%% chuong: 5
%% bai: 16
%% dang: 11.16.1
%% loai: TN4
%% muc: VD
%% dap_an: "A"
%% nguon: "tai-lieu-hoc-tap-mon-toan-lop-11-theo-chuong-trinh-giao-duc-pho-thong-2018.pdf (Lớp Toán Thầy Duy, Phạm Lê Duy, THPT Chu Văn An), Chương 3 Bài 2, Phần 1, Câu 35 (tr.372)"
%% kien_thuc: "Lớp 11 Bài 16 (giới hạn của hàm số)"
%% trang_thai: chinh-thuc
%% ghi_chu: "Đã duyệt 03/11/2026."
\begin{ex}
Biết $\lim\limits_{x\to3}\dfrac{x^2+bx+c}{x-3}=8$ ($b,c\in\mathbb R$). Tính $P=b+c$.
\choice{\True $P=-13$}{$P=-11$}{$P=5$}{$P=-12$}
\loigiai{\sloppy Giới hạn hữu hạn nên tử triệt tiêu tại $x=3$: $9+3b+c=0$, suy ra $c=-9-3b$ và tử $=(x-3)(x+3+b)$. Khi đó giới hạn $=6+b=8\Rightarrow b=2$, $c=-15$. Vậy $P=-13$.}
\end{ex}
